\section{Introduction}
\label{sec:intro}

\IEEEPARstart{L}{arge} language models (LLMs) are now the default writing and translation assistants for many people who work in more than one language. In Kazakhstan, students, translators, journalists and public bodies routinely ask general-purpose LLMs to translate between Kazakh and English or to answer questions about Kazakh text. Kazakh is the state language of the country, but in the vocabulary of language technology it is a low- to mid-resource language: it has comparatively little curated digital text, a Cyrillic script that is shared with Russian, and an agglutinative morphology in which a single word can carry a stem followed by plural, possessive, case and other suffixes chosen by vowel harmony~\cite{joshi2020state,mirzakhalov2021turkic}. Each of these properties is a plausible reason why a model trained mostly on English could perform worse when it must \emph{produce} Kazakh.

Evidence that such a disadvantage exists comes from three directions. First, multilingual benchmarking of generative models reports systematically lower scores for lower-resource languages~\cite{ahuja2023mega,lai2023beyond}, and large translation studies find that general LLMs approach dedicated systems for high-resource languages but lag for low-resource ones~\cite{robinson2023chatgpt,zhu2024multilingual}. Second, Kazakh-specific evaluations reach the same conclusion: KazMMLU shows that even strong models struggle on Kazakh knowledge questions and that accuracy depends heavily on the prompt language~\cite{togmanov2025kazmmlu}, a pilot evaluation of seven LLMs documents weaknesses on Kazakh tasks~\cite{maxutov2024kazakh}, and the KazParC corpus study reports weaker scores for translation into Kazakh than into Russian~\cite{yeshpanov2024kazparc}. Third, the inequality begins before any text is generated: commercial tokenizers need several times more tokens to encode the same content in many non-English languages, which raises cost and shortens the usable context~\cite{petrov2023tokenizer,ahia2023cost}.

These studies, however, leave a specific gap. Each isolates one symptom---subject knowledge, corpus-level translation quality, or token counts---and none measures the \emph{same} items in Kazakh and English for a general-purpose LLM while also asking \emph{which measurable properties of the output} accompany the difference. KazMMLU has no English counterpart, the KazParC evaluation compares dedicated translation systems rather than LLMs, and tokenizer studies measure input length rather than output quality. The explanation most often offered, that agglutinative morphology is to blame, has been suggested but rarely tested~\cite{yeshpanov2024kazparc}. Without paired measurements, a lower Kazakh score cannot be separated from differences in item difficulty, and without diagnostics it cannot inform users or developers about what goes wrong.

This article addresses that gap with a paired, within-item design. Our research question, fixed in the study protocol, is: \emph{How does the accuracy of selected LLMs on Kazakh-language translation and question-answering tasks compare with their accuracy on equivalent English-language tasks, and which types of linguistic error are associated with any differences observed?} We hypothesised that performance is lower in Kazakh (H1) and that morphological errors make up the largest share of Kazakh errors (H2).

The study was designed for three LLMs, a dedicated translation baseline and blind human annotation. At the time of writing, the saved evidence covers one model, Gemini 3.5 Flash-Lite, on \NPairs{} FLORES-200 sentences translated in both directions (\NResponses{} of \NPromptsPerModel{} planned responses for that model). We therefore answer the translation part of the question for this model with full statistical rigour, test morphology-related explanations with automatic, controlled diagnostics rather than with human labels, and report the remaining conditions transparently as pending. All of them can be completed with the released pipeline without changing the protocol.

The contributions of this article are as follows.
\begin{enumerate}
\item \textbf{A paired Kazakh--English benchmark and an auditable pipeline.} KAZLLM contains \NItems{} items (\NPromptsPerModel{} prompts per model) from FLORES-200, KazParC and Belebele, plus 30 team-written probes covering six Kazakh phenomena. Inputs are frozen by SHA-256 hashes, every saved response carries its prompt hash and raw API record, and the scorer refuses incomplete or simulated runs.
\item \textbf{A rigorous directional comparison.} On \NPairs{} paired items, translation into Kazakh is $\GemChrfGapAbs$ chrF2 points worse than translation into English, with a paired bootstrap interval that excludes zero and a consistent direction across four metrics and three sensitivity analyses.
\item \textbf{Controlled diagnostics of the gap.} We introduce an inflection-tolerant prefix analysis with a mismatched-reference control and a truncation ablation, an automatic script audit, and an analysis of API-reported token usage. Together they show that suffix-level mismatch and tokenization cost are measurable contributors, whereas script errors explain only a small part of the gap.
\item \textbf{Transparent reporting.} Every number in the article is generated from the saved logs by one command, and all protocol deviations and pending conditions are documented.
\end{enumerate}

The rest of the article is organised as follows. Section~\ref{sec:related} synthesises related work, Section~\ref{sec:methods} describes the design, data and analysis, Section~\ref{sec:results} reports the results, Section~\ref{sec:discussion} interprets them and discusses threats to validity, and Section~\ref{sec:conclusion} concludes.
